\newpage
% ============================================================
\section*{附录A\quad 文件列表}
\begin{table}[H]
  \centering
  \caption{提交文件列表}
  \label{tab:appendix-files}
  \begin{tabular}{cl}
    \toprule
    文件名 & 功能描述 \\
    \midrule
    q1.py & 问题一程序代码 \\
    q2.py & 问题二程序代码 \\
    q3.py & 问题三程序代码 \\
    q4.py & 问题四程序代码 \\
    \bottomrule
  \end{tabular}
\end{table}

\section*{附录B\quad 部分数据展示}
【占位：展示原始数据或关键中间结果的部分样例，用于说明数据来源与建模输入的真实性。】

\begin{table}[H]
  \centering
  \caption{部分数据展示（示例）}
  \label{tab:appendix-data}
  \begin{tabular}{ccccc}
    \toprule
    编号 & 变量1 & 变量2 & 变量3 & 变量4 \\
    \midrule
    【占位】 & 【占位】 & 【占位】 & 【占位】 & 【占位】 \\
    【占位】 & 【占位】 & 【占位】 & 【占位】 & 【占位】 \\
    \bottomrule
  \end{tabular}
\end{table}

\section*{附录C\quad 代码}

% 方式一：直接内嵌代码（适合较短脚本，或需要展示核心片段时）
\subsection*{q1.py}
\begin{lstlisting}
# -*- coding: utf-8 -*-
"""
【占位：问题一核心代码，请替换为实际脚本内容】
"""
\end{lstlisting}

% 方式二：从外部文件直接引入完整代码（推荐，避免正文与代码文件不同步）
% 取消注释并将路径替换为你们的实际代码路径后使用：
% \lstinputlisting[language=Python, caption={q1.py}]{code/q1.py}
% \lstinputlisting[language=Python, caption={q2.py}]{code/q2.py}
% \lstinputlisting[language=Python, caption={q3.py}]{code/q3.py}
% \lstinputlisting[language=Python, caption={q4.py}]{code/q4.py}
